\chapter{2007年新课标卷（理）}

\section{单选题}

\begin{problem}
已知命题 \(p\)：\(\forall x \in \R\)，\(\sin x \leqslant 1\)，则（\quad）

\choices
    {\(\neg p\)：\(\exists x \in \R\)，\(\sin x \geqslant 1\)}
    {\(\neg p\)：\(\forall x \in \R\)，\(\sin x \geqslant 1\)}
    {\(\neg p\)：\(\exists x \in \R\)，\(\sin x > 1\)}
    {\(\neg p\)：\(\forall x \in \R\)，\(\sin x > 1\)}
\end{problem}

\begin{answer}
C
\end{answer}

\begin{solution}
全称命题的否定要把全称量词改为存在量词，并把不等式 \(\sin x\leqslant 1\) 否定为 \(\sin x>1\). 因此 \(\neg p\) 为“存在 \(x\in\R\)，使 \(\sin x>1\)”，故选 C.
\end{solution}

\begin{problem}
已知平面向量 \(\bs{a} = (1,1)\)，\(\bs{b} = (1,-1)\)，则向量 \(\frac{1}{2}\bs{a} - \frac{3}{2}\bs{b} =\)（\quad）

\choices
    {\((-2, - 1)\)}
    {\((-2, 1)\)}
    {\((-1, 0)\)}
    {\((-1, 2)\)}
\end{problem}

\begin{answer}
D
\end{answer}

\begin{solution}
分别计算两个向量的数乘，得 \(\frac{1}{2}\bs{a}=(\frac{1}{2},\frac{1}{2})\)，\(\frac{3}{2}\bs{b}=(\frac{3}{2},-\frac{3}{2})\). 所以
\(\frac{1}{2}\bs{a}-\frac{3}{2}\bs{b}=(\frac{1}{2}-\frac{3}{2},\frac{1}{2}+\frac{3}{2})=(-1,2)\)，故选 D.
\end{solution}

\begin{problem}
函数 \( y = \sin \left(2x - \frac{\pi}{3}\right) \) 在区间 \(\left[-\frac{\pi}{2}, \pi\right]\) 的简图是（\quad）

\choices
    {\choicebitmap[width=0.15\paperwidth]{img/2007/new_curriculum_science/q3_optA.png}}
    {\choicebitmap[width=0.15\paperwidth]{img/2007/new_curriculum_science/q3_optB.png}}
    {\choicebitmap[width=0.15\paperwidth]{img/2007/new_curriculum_science/q3_optC.png}}
    {\choicebitmap[width=0.15\paperwidth]{img/2007/new_curriculum_science/q3_optD.png}}
\end{problem}

\begin{answer}
A
\end{answer}

\begin{solution}
函数的零点满足 \(2x-\frac{\pi}{3}=k\pi\)，故在给定区间内的零点为 \(-\frac{\pi}{3}\)，\(\frac{\pi}{6}\)，\(\frac{2\pi}{3}\). 又 \(y(-\frac{\pi}{2})=\frac{\sqrt{3}}{2}>0\)，\(y(0)=-\frac{\sqrt{3}}{2}<0\)，\(y(\pi)=-\frac{\sqrt{3}}{2}<0\)，并且函数在 \(-\frac{\pi}{12}\) 处取得极小值 \(-1\)，在 \(\frac{5\pi}{12}\) 处取得极大值 \(1\). 这些特征与选项 A 的图象一致，故选 A.
\end{solution}

\begin{problem}
已知 \(\{a_{n}\}\) 是等差数列，\(a_{10} = 10\)，其前 10 项和 \(S_{10} = 70\)，则其公差 \(d =\)（\quad）

\choices
    {\(-\frac{2}{3}\)}
    {\(-\frac{1}{3}\)}
    {\( \frac{1}{3} \)}
    {\( \frac{2}{3} \)}
\end{problem}

\begin{answer}
D
\end{answer}

\begin{solution}
设等差数列首项为 \(a_1\)，公差为 \(d\). 由 \(S_{10}=\frac{10}{2}(a_1+a_{10})\) 得 \(a_1+10=14\)，所以 \(a_1=4\). 再由 \(a_{10}=a_1+9d=10\)，得 \(4+9d=10\)，从而 \(d=\frac{2}{3}\)，故选 D.
\end{solution}

\begin{problem}
如果执行下面的程序框图，那么输出的 \(S\) =（\quad）

\texfigure[max width=0.56\linewidth]{img_repaint/2007/new_curriculum_science/q5_fig1.tex}

\choices
    {\(2450\)}
    {\(2500\)}
    {\(2550\)}
    {\(2652\)}
\end{problem}

\begin{answer}
C
\end{answer}

\begin{solution}
程序先令 \(k=1\)，\(S=0\)，然后在 \(k\leqslant 50\) 时反复执行 \(S\leftarrow S+2k\)，并令 \(k\leftarrow k+1\). 循环结束时
\(S=2(1+2+\cdots+50)=2\cdot\frac{50\cdot51}{2}=2550\)，故选 C.
\end{solution}

\begin{problem}
已知抛物线 \( y^{2} = 2px \)（\( p > 0 \)）的焦点为 \(F\)，点 \( P_{1}(x_{1}, y_{1}) \)，\( P_{2}(x_{2}, y_{2}) \)，\( P_{3}(x_{3}, y_{3}) \) 在抛物线上，且 \( 2x_{2} = x_{1} + x_{3} \)，则有（\quad）

\choices
    {\( |FP_{1}| + |FP_{2}| = |FP_{3}| \)}
    {\( |FP_{1}|^{2} + |FP_{2}|^{2} = |FP_{3}|^{2} \)}
    {\(2|FP_{2}| = |FP_{1}| + |FP_{3}|\)}
    {\( |FP_{2}|^{2} = |FP_{1}| \cdot |FP_{3}| \)}
\end{problem}

\begin{answer}
C
\end{answer}

\begin{solution}
抛物线 \(y^2=2px\) 的焦点为 \(F(\frac{p}{2},0)\). 对抛物线上任一点 \(P_i(x_i,y_i)\)，由 \(y_i^2=2px_i\) 得
\(\lvert FP_i\rvert^2=(x_i-\frac{p}{2})^2+y_i^2=(x_i+\frac{p}{2})^2\). 由于 \(x_i\geqslant0\)，所以 \(\lvert FP_i\rvert=x_i+\frac{p}{2}\). 于是由 \(2x_2=x_1+x_3\)，有
\(2\lvert FP_2\rvert=2x_2+p=x_1+x_3+p=\lvert FP_1\rvert+\lvert FP_3\rvert\)，故选 C.
\end{solution}

\begin{problem}
已知 \(x > 0\)，\(y > 0\)，\(x\)，\(a\)，\(b\)，\(y\) 成等差数列，\(x\)，\(c\)，\(d\)，\(y\) 成等比数列，则 \(\frac{(a + b)^2}{cd}\) 的最小值是（\quad）

\choices
    {\(0\)}
    {\(1\)}
    {\(2\)}
    {\(4\)}
\end{problem}

\begin{answer}
D
\end{answer}

\begin{solution}
设等差数列的公差为 \(r\)，则 \(a=x+r\)，\(b=x+2r\)，\(y=x+3r\)，从而 \(a+b=x+y\). 设等比数列的公比为 \(q\)，则 \(c=xq\)，\(d=xq^2\)，\(y=xq^3\)，所以 \(cd=x^2q^3=xy\). 因此
\(\frac{(a+b)^2}{cd}=\frac{(x+y)^2}{xy}=\frac{x}{y}+\frac{y}{x}+2\geqslant4\)，等号在 \(x=y\) 时取得，故最小值为 \(4\)，选 D.
\end{solution}

\begin{problem}
已知某个几何体的三视图如下，根据图中标出的尺寸（单位：\(\mathrm{cm}\)），可得这个几何体的体积是（\quad）

\bitmapfigure[max width=0.84\linewidth]{img/2007/new_curriculum_science/q8_fig1.png}

\choices
    {\(\frac{4000}{3}\mathrm{~cm}^{3}\)}
    {\(\frac{8000}{3}\mathrm{~cm}^{3}\)}
    {\(2000\mathrm{~cm}^{3}\)}
    {\(4000\mathrm{~cm}^{3}\)}
\end{problem}

\begin{answer}
B
\end{answer}

\begin{solution}
由三视图可知，该几何体为底面边长为 \(20\mathrm{~cm}\)、高为 \(20\mathrm{~cm}\) 的四棱锥. 其底面为正方形，面积为 \(20^2=400\mathrm{~cm}^2\)，所以体积为
\(V=\frac{1}{3}\cdot400\cdot20=\frac{8000}{3}\mathrm{~cm}^3\)，故选 B.
\end{solution}

\begin{problem}
若 \(\frac{\cos 2\alpha}{\sin \left(\alpha - \frac{\pi}{4}\right)} = -\frac{\sqrt{2}}{2}\)，则 \(\cos \alpha + \sin \alpha\) 的值为（\quad）

\choices
    {\(-\frac{\sqrt{7}}{2}\)}
    {\(-\frac{1}{2}\)}
    {\( \frac{1}{2} \)}
    {\(\frac{\sqrt{7}}{2}\)}
\end{problem}

\begin{answer}
C
\end{answer}

\begin{solution}
由题意分母不为零，故 \(\sin(\alpha-\frac{\pi}{4})\neq0\)，从而 \(\cos\alpha-\sin\alpha\neq0\). 利用
\(\cos2\alpha=(\cos\alpha+\sin\alpha)(\cos\alpha-\sin\alpha)\)，以及 \(\sin(\alpha-\frac{\pi}{4})=\frac{\sin\alpha-\cos\alpha}{\sqrt{2}}\)，原式化为
\(-\sqrt{2}(\cos\alpha+\sin\alpha)=-\frac{\sqrt{2}}{2}\). 所以 \(\cos\alpha+\sin\alpha=\frac{1}{2}\)，故选 C.
\end{solution}

\begin{problem}
曲线 \( y = \e^{\frac{1}{2} x} \) 在点 \((4, \e^2)\) 处的切线与坐标轴所围三角形的面积为（\quad）

\choices
    {\(\frac{9}{2}\e^2\)}
    {\(4\e^2\)}
    {\(2\e^{2}\)}
    {\(\e^2\)}
\end{problem}

\begin{answer}
D
\end{answer}

\begin{solution}
函数在 \((4,\e^2)\) 处的导数为 \(y'=\frac{1}{2}\e^{x/2}\)，所以切线斜率为 \(\frac{1}{2}\e^2\). 切线方程为
\(y-\e^2=\frac{1}{2}\e^2(x-4)\)，即 \(y=\frac{1}{2}\e^2x-\e^2\). 它与 \(x\) 轴交于 \((2,0)\)，与 \(y\) 轴交于 \((0,-\e^2)\)，故所围三角形面积为 \(\frac{1}{2}\cdot2\cdot\e^2=\e^2\)，选 D.
\end{solution}

\begin{problem}
甲，乙，丙三名射箭运动员在某次测试中各射箭 \(20\) 次，三人的测试成绩如下表

\begin{center}
\begin{tabular}{ccccc}
\multicolumn{5}{c}{甲的成绩} \\
环数 & \(7\) & \(8\) & \(9\) & \(10\) \\
频数 & \(5\) & \(5\) & \(5\) & \(5\)
\end{tabular}
\qquad
\begin{tabular}{ccccc}
\multicolumn{5}{c}{乙的成绩} \\
环数 & \(7\) & \(8\) & \(9\) & \(10\) \\
频数 & \(6\) & \(4\) & \(4\) & \(6\)
\end{tabular}
\qquad
\begin{tabular}{ccccc}
\multicolumn{5}{c}{丙的成绩} \\
环数 & \(7\) & \(8\) & \(9\) & \(10\) \\
频数 & \(4\) & \(6\) & \(6\) & \(4\)
\end{tabular}
\end{center}

\(s_{1}\)，\(s_{2}\)，\(s_{3}\) 分别表示甲，乙，丙三名运动员这次测试成绩的标准差，则有（\quad）

\choices
    {\(s_3 > s_1 > s_2\)}
    {\(s_2 > s_1 > s_3\)}
    {\(s_1 > s_2 > s_3\)}
    {\(s_2 > s_3 > s_1\)}
\end{problem}

\begin{answer}
B
\end{answer}

\begin{solution}
三名运动员的成绩平均数均为 \(8.5\). 按标准差公式计算平方和，甲、乙、丙的方差分别为
\(s_1^2=\frac{5(1.5^2+0.5^2+0.5^2+1.5^2)}{20}=\frac{25}{20}\)，
\(s_2^2=\frac{6\cdot1.5^2+4\cdot0.5^2+4\cdot0.5^2+6\cdot1.5^2}{20}=\frac{29}{20}\)，
\(s_3^2=\frac{4\cdot1.5^2+6\cdot0.5^2+6\cdot0.5^2+4\cdot1.5^2}{20}=\frac{21}{20}\). 因此 \(s_2>s_1>s_3\)，故选 B.
\end{solution}

\begin{problem}
一个四棱锥和一个三棱锥恰好可以拼接成一个三棱柱，这个四棱锥的底面为正方形，且底面边长与各侧棱长相等，这个三棱锥的底面边长与各侧棱长也都相等. 设四棱锥，三棱锥，三棱柱的高分别为 \(h_1\)，\(h_2\)，\(h\) 则 \(h_1: h_2: h =\)（\quad）

\choices
    {\(\sqrt{3}: 1: 1\)}
    {\(\sqrt{3}: 2: 2\)}
    {\(\sqrt{3}: 2: \sqrt{2}\)}
    {\(\sqrt{3}: 2: \sqrt{3}\)}
\end{problem}

\begin{answer}
B
\end{answer}

\begin{solution}
设拼接时各相关棱长为 \(a\). 四棱锥的底面是边长为 \(a\) 的正方形，顶点到底面中心的距离为 \(a/\sqrt{2}\)，因此
\(h_1=\sqrt{a^2-(a/\sqrt{2})^2}=\frac{a}{\sqrt{2}}\). 三棱锥的底面是边长为 \(a\) 的正三角形，底面中心到顶点的距离为 \(a/\sqrt{3}\)，所以
\(h_2=\sqrt{a^2-(a/\sqrt{3})^2}=a\sqrt{\frac{2}{3}}\).

拼接后的三棱柱以该正三角形为底面，底面积为 \(\frac{\sqrt{3}}{4}a^2\). 由体积可加性，且四棱锥底面面积为 \(a^2\)，有
\(\frac{\sqrt{3}}{4}a^2h=\frac{1}{3}a^2h_1+\frac{1}{3}\cdot\frac{\sqrt{3}}{4}a^2h_2\). 代入 \(h_1\)、\(h_2\) 得 \(h=a\sqrt{\frac{2}{3}}=h_2\). 因此
\(h_1:h_2:h=\frac{a}{\sqrt{2}}:a\sqrt{\frac{2}{3}}:a\sqrt{\frac{2}{3}}=\sqrt{3}:2:2\)，故选 B.
\end{solution}

\section{填空题}

\begin{problem}
已知双曲线的顶点到渐近线的距离为 \(2\)，焦点到渐近线的距离为 \(6\)，则该双曲线的离心率为\fillinblank{}.
\end{problem}

\begin{answer}
\(3\)
\end{answer}

\begin{solution}
设双曲线为 \(\frac{x^2}{a^2}-\frac{y^2}{b^2}=1\)，焦距参数满足 \(c^2=a^2+b^2\). 取渐近线 \(y=\frac{b}{a}x\)，写成 \(bx-ay=0\). 顶点 \((a,0)\) 到该渐近线的距离为 \(\frac{ab}{c}=2\)，焦点 \((c,0)\) 到该渐近线的距离为 \(b=6\). 于是 \(\frac{a}{c}=\frac{1}{3}\)，所以离心率 \(e=\frac{c}{a}=3\).
\end{solution}

\begin{problem}
设函数 \( f(x) = \frac{(x + 1)(x + a)}{x} \) 为奇函数，则 \( a = \)\fillinblank{}.
\end{problem}

\begin{answer}
\(-1\)
\end{answer}

\begin{solution}
\(f(x)=\frac{(x+1)(x+a)}{x}=x+(a+1)+\frac{a}{x}\)，定义域为 \(x\neq0\). 奇函数要求 \(f(-x)=-f(x)\). 比较两边可得
\(-x+(a+1)-\frac{a}{x}=-x-(a+1)-\frac{a}{x}\)，所以 \(a+1=-(a+1)\)，即 \(a=-1\). 此时 \(f(x)=x-\frac{1}{x}\)，确为奇函数.
\end{solution}

\begin{problem}
\(\i\) 是虚数单位，\( \frac{-5 + 10\i}{3 + 4\i} = \)\fillinblank{}. （用 \(a + b\i\) 的形式表示，\(a\)，\(b \in \R\)）
\end{problem}

\begin{answer}
\(1+2\i\)
\end{answer}

\begin{solution}
分子分母同乘分母的共轭复数 \(3-4\i\)，得
\(\frac{-5+10\i}{3+4\i}=\frac{(-5+10\i)(3-4\i)}{(3+4\i)(3-4\i)}=\frac{25+50\i}{25}=1+2\i\).
\end{solution}

\begin{problem}
某校安排 \(5\) 个班到 \(4\) 个工厂进行社会实践. 每个班去一个工厂，每个工厂至少安排一个班，不同的安排方法共有\fillinblank{} 种. （用数字作答）
\end{problem}

\begin{answer}
\(240\)
\end{answer}

\begin{solution}
为了使四个工厂均有班级，5 个有标号的班级中必有一个工厂分配到两个班级，其余三个工厂各分配一个班级. 先选出分配两个班级的工厂，有 \(4\) 种；再从5个班级中选出这两个班级，有 \(\mathrm C_5^2\) 种；最后将剩下的三个班级分配给其余三个工厂，有 \(3!\) 种. 因此安排方法数为
\(4\mathrm C_5^2\cdot3!=4\cdot10\cdot6=240\).
\end{solution}

\section{解答题}

\begin{problem}
如图，测量河对岸的塔高 \(AB\) 时，可以选择与塔底 \(B\) 在同一水平面内的两个测点 \(C\) 与 \(D\). 现测得 \(\angle BCD = \alpha\)，\(\angle BDC = \beta\)，\(CD = s\)，并在点 \(C\) 测得塔顶 \(A\) 的仰角为 \(\theta\)，求塔高 \(AB\).

\FigureLayoutDeclare{img/2007/new_curriculum_science/q17_fig1.png}{10.5cm}{108bp 65bp 116bp 47bp}
\bitmapfigure[width=10.5cm]{img/2007/new_curriculum_science/q17_fig1.png}
\end{problem}

\begin{solution}
在 \(\triangle BCD\) 中，\(\angle CBD=\pi-\alpha-\beta\). 由正弦定理，
\(\frac{BC}{\sin\beta}=\frac{CD}{\sin(\pi-\alpha-\beta)}\)，所以
\(BC=\frac{s\sin\beta}{\sin(\alpha+\beta)}\).

由于塔高 \(AB\) 垂直于测点所在的水平面，\(\triangle ABC\) 是直角三角形，且 \(\angle ACB=\theta\). 因此
\(AB=BC\tan\theta=\frac{s\sin\beta\tan\theta}{\sin(\alpha+\beta)}\).
\end{solution}

\begin{problem}
如图，在三棱锥 \(S - ABC\) 中，侧面 \(SAB\) 与侧面 \(SAC\) 均为等边三角形，\(\angle BAC = 90^{\circ}\)，\(O\) 为 \(BC\) 中点.

\begin{enumerate}
\item 证明：\(SO \perp\) 平面 \(ABC\)；
\item 求二面角 \(A - SC - B\) 的余弦值.
\end{enumerate}

\bitmapfigure[max width=0.48\linewidth]{img/2007/new_curriculum_science/q18_fig1.png}
\end{problem}

\begin{solution}
\begin{enumerate}
\item
设等边三角形的边长为 \(a\)，则 \(AB=AC=SA=SB=SC=a\). 因为 \(O\) 是 \(BC\) 的中点，且 \(AB=AC\)，所以在等腰三角形 \(ABC\) 中 \(AO\perp BC\)；在等腰三角形 \(SBC\) 中 \(SO\perp BC\). 又
\(AO=\frac{a}{\sqrt{2}}\)，\(OB=\frac{a}{\sqrt{2}}\)，从 \(\triangle SOB\) 得 \(SO=\frac{a}{\sqrt{2}}\). 于是 \(SO^2+AO^2=a^2=SA^2\)，故 \(SO\perp AO\). 直线 \(AO\) 与 \(BC\) 是平面 \(ABC\) 内相交的两条直线，所以 \(SO\perp\) 平面 \(ABC\).

\item
取 \(SC\) 的中点 \(M\). 在等边三角形 \(SAC\) 中，\(AM\perp SC\). 又 \(O\)，\(M\) 分别是 \(BC\)，\(SC\) 的中点，所以 \(OM\parallel SB\). 由 \(BC=\sqrt{AB^2+AC^2}=a\sqrt{2}\) 及 \(SB=SC=a\)，得 \(\angle BSC=90^{\circ}\)，从而 \(OM\perp SC\). 因此 \(\angle AMO\) 是二面角 \(A-SC-B\) 的平面角. 由上面的垂直关系，\(AO\perp SO\) 且 \(AO\perp BC\)，而 \(SO\) 与 \(BC\) 是平面 \(SBC\) 内相交的两条直线，故 \(AO\perp\) 平面 \(SBC\)，从而 \(AO\perp OM\). 在直角三角形 \(AOM\) 中，\(OM=\frac{a}{2}\)，\(AO=\frac{a}{\sqrt{2}}\)，所以 \(AM=\frac{\sqrt{3}}{2}a\)，从而
\(\cos\angle AMO=\frac{OM}{AM}=\frac{1}{\sqrt{3}}\).
\end{enumerate}
\end{solution}

\begin{problem}
在平面直角坐标系 \(xOy\) 中，经过点 \((0, \sqrt{2})\) 且斜率为 \(k\) 的直线 \(l\) 与椭圆 \(\frac{x^2}{2} + y^2 = 1\) 有两个不同的交点 \(P\) 和 \(Q\).

\begin{enumerate}
\item 求 \(k\) 的取值范围；
\item 设椭圆与 \(x\) 轴正半轴，\(y\) 轴正半轴的交点分别为 \(A\)，\(B\)，是否存在常数 \(k\)，使得向量 \(\overrightarrow{OP} + \overrightarrow{OQ}\) 与 \(\overrightarrow{AB}\) 共线？如果存在，求 \(k\) 值；如果不存在，请说明理由.
\end{enumerate}
\end{problem}

\begin{solution}
\begin{enumerate}
\item 直线 \(l\) 经过 \((0,\sqrt{2})\)，所以方程为 \(y=kx+\sqrt{2}\). 代入椭圆方程，得
\(\frac{x^2}{2}+(kx+\sqrt{2})^2=1\)，即 \(\left(k^2+\frac{1}{2}\right)x^2+2\sqrt{2}kx+1=0\).

直线与椭圆有两个不同交点，当且仅当上述关于 \(x\) 的二次方程有两个不等实根，即
\(\Delta=(2\sqrt{2}k)^2-4\left(k^2+\frac{1}{2}\right)=4k^2-2>0\). 因此 \(k< -\frac{\sqrt{2}}{2}\) 或 \(k>\frac{\sqrt{2}}{2}\)，即
\(k\in\left(-\infty,-\frac{\sqrt{2}}{2}\right)\cup\left(\frac{\sqrt{2}}{2},+\infty\right)\).
\item 设 \(P(x_1,y_1)\)、\(Q(x_2,y_2)\). 由上面二次方程的韦达定理，
\(x_1+x_2=-\frac{2\sqrt{2}k}{k^2+\frac{1}{2}}=-\frac{4\sqrt{2}k}{2k^2+1}\)，并且 \(y_1+y_2=k(x_1+x_2)+2\sqrt{2}\). 椭圆与坐标轴正半轴的交点为 \(A(\sqrt{2},0)\)、\(B(0,1)\)，所以 \(\overrightarrow{AB}=(-\sqrt{2},1)\).

由 \(y_i=kx_i+\sqrt{2}\)，有 \(y_1+y_2=k(x_1+x_2)+2\sqrt{2}=\frac{2\sqrt{2}}{2k^2+1}\). 因此向量 \(\overrightarrow{OP}+\overrightarrow{OQ}=(x_1+x_2,y_1+y_2)\) 与 \(\overrightarrow{AB}\) 共线的条件是
\((x_1+x_2)+\sqrt{2}(y_1+y_2)=0\). 代入上式，得
\(-\frac{4\sqrt{2}k}{2k^2+1}+\frac{4}{2k^2+1}=0\). 由于 \(2k^2+1>0\)，所以 \(1-\sqrt{2}k=0\)，即 \(k=\frac{\sqrt{2}}{2}\). 但这不满足第一问的严格不等式，故不存在符合条件的常数 \(k\).
\end{enumerate}
\end{solution}

\begin{problem}
如图，面积为 \(S\) 的正方形 \(ABCD\) 中有一个不规则的图形 \(M\)，可按下面方法估计 \(M\) 的面积：在正方形 \(ABCD\) 中随机投掷 \(n\) 个点，若 \(n\) 个点中有 \(m\) 个点落入 \(M\) 中，则 \(M\) 的面积的估计值为 \(\frac{m}{n} S\). 假设正方形 \(ABCD\) 的边长为 \(2\)，\(M\) 的面积为 \(1\)，并向正方形 \(ABCD\) 中随机投掷 \(10000\) 个点，以 \(X\) 表示落入 \(M\) 中的点的数目.

\begin{enumerate}
\item 求 \(X\) 的均值 \(E(X)\)；
\item 求用以上方法估计 \(M\) 的面积时，\(M\) 的面积的估计值与实际值之差在区间 \((-0.03, 0.03)\) 内的概率. 附表：\(P(k)=\sum_{i=0}^{k}\mathrm C_{10000}^{i}\times 0.25^{i}\times 0.75^{10000-i}\).
\end{enumerate}

\begin{center}
\begin{tabular}{|c|c|c|c|c|}
\hline
\(k\) & \(2424\) & \(2425\) & \(2574\) & \(2575\) \\
\hline
\(P(k)\) & \(0.0403\) & \(0.0423\) & \(0.9570\) & \(0.9590\) \\
\hline
\end{tabular}
\end{center}

\bitmapfigure[max width=0.34\linewidth]{img/2007/new_curriculum_science/q20_fig1.png}
\end{problem}

\begin{solution}
\begin{enumerate}
\item 每个随机点落入 \(M\) 的概率为 \(\frac{1}{4}\)，所以 \(X\sim B\left(10000,\frac{1}{4}\right)\). 因此
\(E(X)=10000\cdot\frac{1}{4}=2500\).
\item 估计面积为 \(\frac{4X}{10000}=\frac{X}{2500}\). 题设不等式等价于
\(-0.03<\frac{X}{2500}-1<0.03\)，即 \(2425<X<2575\). 由于 \(X\) 为整数，所求概率为
\(P(2426\leqslant X\leqslant2574)=P(2574)-P(2425)=0.9570-0.0423=0.9147\).
\end{enumerate}
\end{solution}

\begin{problem}
设函数 \( f(x) = \ln (x + a) + x^2 \).

\begin{enumerate}
\item 若当 \(x = -1\) 时，\(f(x)\) 取得极值，求 \(a\) 的值，并讨论 \(f(x)\) 的单调性；
\item 若 \(f(x)\) 存在极值，求 \(a\) 的取值范围，并证明所有极值之和大于 \( \ln \frac{\e}{2} \).
\end{enumerate}
\end{problem}

\begin{solution}
\begin{enumerate}
\item 函数定义域为 \(x>-a\)，且 \(f'(x)=\frac{1}{x+a}+2x\). 由于 \(x=-1\) 是定义域内的极值点，必须有 \(f'(-1)=0\)，即 \(\frac{1}{a-1}-2=0\)，所以 \(a=\frac{3}{2}\). 此时
\(f'(x)=\frac{2x^2+3x+1}{x+\frac{3}{2}}=\frac{(2x+1)(x+1)}{x+\frac{3}{2}}\)，定义域为 \(\left(-\frac{3}{2},+\infty\right)\). 据此，\(f'(x)>0\) 在 \(\left(-\frac{3}{2},-1\right)\) 和 \(\left(-\frac{1}{2},+\infty\right)\) 上成立，\(f'(x)<0\) 在 \(\left(-1,-\frac{1}{2}\right)\) 上成立. 因此函数在 \(\left(-\frac{3}{2},-1\right)\) 和 \(\left(-\frac{1}{2},+\infty\right)\) 上单调递增，在 \(\left(-1,-\frac{1}{2}\right)\) 上单调递减.
\item 令 \(t=x+a\)，则定义域为 \(t>0\)，且
\(f'(x)=2t+\frac{1}{t}-2a\). 函数 \(2t+\frac{1}{t}\) 在 \(t>0\) 上的最小值为 \(2\sqrt{2}\). 当 \(a\leqslant\sqrt{2}\) 时，\(f'(x)\) 没有两个变号零点；当 \(a=\sqrt{2}\) 时只有一个不变号的零点；当 \(a>\sqrt{2}\) 时有两个正根，且导数符号依次为正、负、正，函数存在极大值和极小值. 因此存在极值的充要条件为 \(a>\sqrt{2}\).

当 \(a>\sqrt{2}\) 时，设两个临界点对应的 \(t\) 值为 \(t_1\)，\(t_2\). 由 \(2t^2-2at+1=0\)，有 \(t_1+t_2=a\)，\(t_1t_2=\frac{1}{2}\). 相应的 \(x_i=t_i-a\)，两个极值之和为
\(f(x_1)+f(x_2)=\ln(t_1t_2)+(t_1-a)^2+(t_2-a)^2=a^2-1-\ln2\). 因为 \(a^2>2\)，所以
\(f(x_1)+f(x_2)>1-\ln2=\ln\frac{\e}{2}\).
\end{enumerate}
\end{solution}

\begin{problem}
如图，已知 \(AP\) 是 \(\odot O\) 的切线，\(P\) 为切点，直线 \(AC\) 是 \(\odot O\) 的割线，与 \(\odot O\) 交于 \(B\)，\(C\) 两点，圆心 \(O\) 在 \(\angle PAC\) 的内部，点 \(M\) 是 \(BC\) 的中点.

\begin{enumerate}
\item 证明：\(A\)，\(P\)，\(O\)，\(M\) 四点共圆；
\item 求 \( \angle OAM + \angle APM \) 的大小.
\end{enumerate}

\bitmapfigure[max width=0.48\linewidth]{img/2007/new_curriculum_science/q22_fig1.png}
\end{problem}

\begin{solution}
\begin{enumerate}
\item
连接 \(OP\)、\(OM\). 因为 \(AP\) 是圆在 \(P\) 处的切线，所以 \(OP\perp AP\)；又因为 \(M\) 是弦 \(BC\) 的中点，所以 \(OM\perp BC\). 由于 \(A\)、\(B\)、\(C\) 共线，\(AM\) 与 \(BC\) 共线，从而 \(OM\perp AM\). 因此 \(\angle OPA=\angle OMA=90^{\circ}\)，四边形 \(APOM\) 的一组对角互补，故 \(A\)、\(P\)、\(O\)、\(M\) 四点共圆.

\item
由四点共圆，\(\angle OAM=\angle OPM\). 又因 \(OP\perp AP\)，且圆心 \(O\) 在 \(\angle PAC\) 内部，射线 \(PM\) 位于由射线 \(PA\)、\(PO\) 构成的直角内，所以 \(\angle OPM+\angle APM=90^{\circ}\). 于是
\(\angle OAM+\angle APM=90^{\circ}\).
\end{enumerate}
\end{solution}

\section{选考题：解答题}

\begin{problem}
\(\odot O_{1}\) 和 \(\odot O_{2}\) 的极坐标方程分别为 \(\rho = 4\cos \theta\)，\(\rho = -4\sin \theta\).

\begin{enumerate}
\item 把 \(\odot O_{1}\) 和 \(\odot O_{2}\) 的极坐标方程化为直角坐标方程；
\item 求经过 \(\odot O_{1}\)，\(\odot O_{2}\) 交点的直线的直角坐标方程.
\end{enumerate}
\end{problem}

\begin{solution}
\begin{enumerate}
\item 在极坐标与直角坐标的同一原点、同一极轴下，\(x=\rho\cos\theta\)，\(y=\rho\sin\theta\)，\(\rho^2=x^2+y^2\). 由 \(\rho=4\cos\theta\)，两边乘以 \(\rho\) 得 \(\rho^2=4x\)，所以 \(\odot O_1\) 的直角坐标方程为 \(x^2+y^2-4x=0\). 对 \(\rho=-4\sin\theta\)，两边乘以 \(\rho\) 得 \(\rho^2=-4y\)，所以 \(\odot O_2\) 的直角坐标方程为 \(x^2+y^2+4y=0\).
\item 两圆的交点满足
\(\begin{cases}
x^2+y^2-4x=0,\\
x^2+y^2+4y=0.
\end{cases}\)
两式相减得 \(x+y=0\)，即 \(y=-x\). 代入任一圆可得交点为 \((0,0)\) 和 \((2,-2)\)，所以经过两个交点的直线方程为 \(y=-x\).
\end{enumerate}
\end{solution}

\begin{problem}
设函数 \( f(x) = |2x + 1| - |x - 4| \).

\begin{enumerate}
\item 解不等式 \( f(x) > 2 \)；
\item 求函数 \( y = f(x) \) 的最小值.
\end{enumerate}
\end{problem}

\begin{solution}
\begin{enumerate}
\item 按 \(2x+1\) 和 \(x-4\) 的符号分段，得
\(f(x)=\begin{cases}
-x-5,&x<-\frac{1}{2},\\
3x-3,&-\frac{1}{2}\leqslant x<4,\\
x+5,&x\geqslant4.
\end{cases}\)
当 \(x<-\frac{1}{2}\) 时，\(-x-5>2\) 等价于 \(x<-7\)；当 \(-\frac{1}{2}\leqslant x<4\) 时，\(3x-3>2\) 等价于 \(x>\frac{5}{3}\)；当 \(x\geqslant4\) 时，\(x+5>2\) 恒成立. 因此不等式的解集为 \(\left(-\infty,-7\right)\cup\left(\frac{5}{3},+\infty\right)\).
\item 第一段函数递减，第二段和第三段函数递增，所以最小值在分界点 \(x=-\frac{1}{2}\) 处取得，且
\(f(-\frac{1}{2})=\left|0\right|-\left|-\frac{9}{2}\right|=-\frac{9}{2}\).
\end{enumerate}
\end{solution}
